\section{No-Code Integration: The \texttt{dashai-frankenstein} Plugin}
\label{sec:dashai-plugin}

The schema-driven CLI and web builder described in Section~\ref{sec:system-design} serve researchers who are comfortable with YAML configuration files. A complementary audience---domain specialists, educators, and analysts who need a trained model but not a training script---is better served by a graphical, no-code environment. To reach that audience, Frankenstein Transformer ships a companion adapter package, \texttt{dashai-frankenstein}, that registers the Frankenstein model classes as components of \textbf{dashAI} (\texttt{\url{https://github.com/DashAISoftware/dashAI}}), an open-source, local-first desktop workbench for machine learning developed by a Chilean academic consortium (Universidad de Chile/FCFM, CENIA, IMFD, and ANID funding). Because both systems are schema-driven at heart---dashAI renders its configuration forms directly from Pydantic schemas, while Frankenstein validates its YAML against a strict JSON Schema---the integration required no changes to dashAI's backend: the plugin is discovered through the standard \texttt{dashai.plugins} entry-points group on startup, exactly as documented in dashAI's plugin guide (\texttt{\url{https://docs.dash-ai.com/learn/guides/plugins/}}), where plugins are ordinary PyPI packages whose first-party generative capabilities are themselves distributed as plugins.

\subsection{Registered Components}

A single \texttt{dashai.plugins} entry-points group registers five components in the dashAI component registry, covering the full breadth of the toolkit's model classes:

\begin{itemize}[leftmargin=*]
    \item \textbf{\texttt{FrankensteinMLMModel}} (\texttt{BaseModel}) --- binds to dashAI's \texttt{TextClassificationTask}. It trains the Frankenstein encoder with a full-precision classification head pooled over the encoder output (Strategy~A of the integration audit), inheriting all sequence mixers, normalizers, and optimizer families from the underlying engine; \texttt{predict} returns a per-class probability matrix.
    \item \textbf{\texttt{FrankensteinDecoderModel}} (\texttt{BaseGenerativeModel}) --- binds to \texttt{TextToTextGenerationTask}. It forces \texttt{model\_class: frankensteindecoder} (which in turn forces \texttt{mode: decoder}) and exposes autoregressive top-\emph{k} sampling; generation-specific parameters (\texttt{max\_new\_tokens}, \texttt{temperature}, \texttt{top\_k}) remain as native dashAI form fields because they are inference-time knobs with no home in the Frankenstein training schema.
    \item \textbf{\texttt{FrankensteinViTClassifier}} (\texttt{BaseModel}) --- binds to \texttt{ImageClassificationTask}, reusing the Vision Transformer's built-in classification head (Appendix~\ref{sec:annex-vision}).
    \item \textbf{\texttt{FrankensteinViTSegmenter}} (\texttt{BaseModel}) --- binds to a plugin-provided \texttt{SegmentationTask}, a new \texttt{BaseTask} component that dashAI does not ship; it uses the ViT segmentation head and returns per-pixel class maps.
    \item \textbf{\texttt{SegmentationTask}} (\texttt{BaseTask}) --- the companion task declaring image-in/mask-out column types for the segmenter.
\end{itemize}

All five components share one thin engine facade over Frankenstein's reusable engine API (the non-CLI, non-supervisor interface extracted in Phase~0 of the integration plan), so there is exactly one heavy dependency (\texttt{frankenstein-transformer}) behind the plugin.

\subsection{Passthrough Configuration Schema}

The integration preserves Frankenstein's schema-first philosophy (Section~\ref{sec:system-design}) by keeping the Frankenstein JSON Schema as the single source of truth rather than translating its $150+$ model and $38$ training fields into native dashAI form fields---a translation that would be brittle against an actively evolving schema. Instead, each component exposes a minimal Pydantic schema whose only user-facing field, \texttt{frankenstein\_json}, is a string containing a complete Frankenstein training config as a \emph{single-line JSON} document (a concession to dashAI's single-line text fields; users build the YAML with the schema-driven web builder and convert it with a one-liner). Before any training or model construction launches, the adapter runs a three-stage validation pipeline:

\begin{enumerate}[leftmargin=*]
    \item \textbf{JSON parsing} --- catches syntax errors in the single-line payload.
    \item \textbf{JSON Schema validation} --- \texttt{jsonschema.validate} against the resolved Frankenstein schema, enforcing \texttt{additionalProperties: false} and every enum range (mixers, optimizers, norms, activations).
    \item \textbf{Config-loader validation} --- Frankenstein's own \texttt{load\_training\_config} enforces the cross-component constraints (divisibility of \texttt{hidden\_size} by \texttt{num\_heads}, \texttt{num\_kv\_heads} dividing \texttt{num\_heads}, vocabulary/tokenizer agreement, BitNet flag rules, optimizer presence per task, \texttt{frankensteindecoder} forcing \texttt{mode: decoder}).
\end{enumerate}

Any failure surfaces to the dashAI user as a single readable \texttt{ValueError} tagged with the failing stage, so invalid configurations never reach the training loop---the same fail-fast contract the CLI provides. Dataset injection is handled by adapters that map dashAI's \texttt{DashAIDataset} (a Hugging Face \texttt{datasets} wrapper) onto Frankenstein's dataset expectations, inject \texttt{vocab\_size} from the resolved tokenizer so the embedding matrix matches the vocabulary, and stream per-epoch metrics back into dashAI's metric store so training progress is visible in the UI.

\subsection{Work in Progress}

The plugin is a work in progress. Its phases are implemented---the MLM text-classification MVP (Phase~1), the generative decoder and ViT classifier (Phase~2), and the ViT segmenter plus the plugin-provided \texttt{SegmentationTask} (Phase~3)---but the package has not yet been executed inside a live dashAI instance: the component contract was implemented from an audit of dashAI's plugin registry, base classes, and job-queue architecture rather than validated end-to-end. Publication to PyPI is likewise pending, both for the plugin itself and for its dependency floor (\texttt{frankenstein-transformer>=1.1.0}, the engine-API release that supersedes the pre-engine \texttt{1.0.0} already on PyPI). Known limitations carried into the next iteration include: the MLM component runs a direct supervised classification loop rather than an MLM-pretrain-then-finetune recipe; the passthrough schema defers curated native form fields for high-leverage knobs (model class, mixer pattern, optimizer family) to a future version; and dashAI's ability to render segmentation-mask output columns is unconfirmed and may require a local frontend addition. These items are tracked in the repository's integration audit (\texttt{\url{https://github.com/erickfmm/frankenstein-transformer}}), and the design intentionally leaves dashAI's backend core untouched, so future dashAI releases cannot break the adapter through coupling.

The integration completes the toolkit's accessibility story: the same validated configuration contract that drives the CLI, the Streamlit builder, and the export pipeline now also drives a third-party no-code workbench, allowing Frankenstein's architecture space to be explored from a graphical interface without sacrificing the schema guarantees on which reproducibility depends.